% Part: incompleteness
% Chapter: theories-computability
% Section: complete-decidable

\documentclass[../../../include/open-logic-section]{subfiles}

\begin{document}

\olfileid{inc}{tcp}{cdc}

\olsection{সম্পূর্ণ \printtoken{S}{axiomatizable} তত্ত্বগুলি নির্ণেয়}

কোনো তত্ত্বকে \emph{সম্পূর্ণ} বলা হয় যদি প্রতিটি বাক্য~$!A$-এর জন্য
$!A$ অথবা $\lnot !A$ প্রমাণযোগ্য হয়।

\begin{lem}
ধরি কোনো তত্ত্ব~$\Th{T}$ সম্পূর্ণ এবং !!{axiomatizable}। তাহলে
$\Th{T}$ নির্ণেয়।
\end{lem}

\begin{proof}
ধরি $\Th{T}$ সম্পূর্ণ এবং $A$ তার একটি গণনসাধ্য স্বতঃসিদ্ধসমষ্টি।
$\Th{T}$ অসঙ্গতিপূর্ণ হলে সেটি স্পষ্টতই গণনসাধ্য। (অ্যালগরিদম: ``শুধু
হ্যাঁ বলো।'') তাই আমরা ধরে নিতে পারি যে $\Th{T}$ সঙ্গতিপূর্ণও।

কোনো বাক্য~$!A$ $\Th{T}$-তে আছে কি না নির্ধারণ করতে, গণনসাধ্য
স্বতঃসিদ্ধসমষ্টি~$A$ থেকে $!A$-এর কোনো !!a{derivation} এবং $A$ থেকে
$\lnot !A$-এর কোনো !!a{derivation} একই সঙ্গে খুঁজতে থাকো। $\Th{T}$
সম্পূর্ণ বলে এ দুটির একটি অবশ্যই পাওয়া যাবে; আর $\Th{T}$ সঙ্গতিপূর্ণ
বলে $\lnot !A$-এর কোনো !!a{derivation} পাওয়া গেলে $!A$-এর কোনো
!!{derivation} নেই।
% BN-SRC-239 TARGET-CORRECTION-BEGIN
\par\smallskip\noindent\textnormal{\small\textbf{উৎস-সংশোধন (BN-SRC-239):}
হিমায়িত প্রমাণে প্রথম অনুসন্ধানটি $\Th{T}$ থেকে এবং দ্বিতীয়টি কোনো
উৎস উল্লেখ না করে চালাতে বলা হয়েছে। কিন্তু কার্যকরভাবে দেওয়া বস্তুটি
হলো গণনসাধ্য স্বতঃসিদ্ধসমষ্টি~$A$; তত্ত্বের সব উপপাদ্যকে সরাসরি
স্বতঃসিদ্ধ ধরে অনুসন্ধান করা চক্রাকার হবে। বাংলা পাঠে উভয় নিষ্পাদনই
$A$ থেকে অনুসন্ধান করা হয়েছে।}
% BN-SRC-239 TARGET-CORRECTION-END

কথাটি অন্যভাবে বলা যায়। আমরা ইতিমধ্যে জানি যে $\Th{T}$ !!{c.e.}; তাই
আগে প্রমাণিত একটি উপপাদ্য অনুসারে $\Th{T}$-এর পূরকটিও !!{c.e.} দেখালেই
যথেষ্ট। কিন্তু কোনো !!a{formula}~$!A$ ঠিক তখনই $\Th{\bar T}$-তে থাকে
যখন $\lnot !A$ $\Th{T}$-তে থাকে; তাই $\Th{\bar T} \leq_m \Th{T}$।
\end{proof}

\end{document}
